% !TeX root = ../main.tex
\chapter{Testing and Evaluation}
\label{chap:testing_evaluation}

\section{Introduction}

This chapter evaluates the functional fidelity, operational resilience, and throughput performance of the proposed predictive maintenance and anomaly detection platform. The evaluation assesses both the physical accuracy of the generated telemetry and the stability of the distributed streaming lakehouse under sustained, real-time industrial workloads. \textit{(To complete: additional evaluation sections and machine learning validation procedures will be incorporated as the platform matures.)}

\paragraph{Experimental Test Environment and Reproducibility.}
To guarantee strict environment parity, reproducibility, and process isolation, the entire data platform is deployed as a containerized stack orchestrated via Docker Compose. Every functional component---including the native Rust-based vehicle simulator, the Eclipse Mosquitto MQTT broker, the Telegraf ingestion bridge, the Redpanda distributed streaming backbone, the Apache Spark Structured Streaming engines, the Apache Iceberg lakehouse on MinIO S3, the QuestDB time-series database, and the Grafana visualization service---executes within independent container boundaries connected through dedicated software-defined bridge networks.

All experimental tests and soak evaluations reported in this chapter were carried out on a dedicated physical workstation. Table~\ref{tab:test-hardware-environment} summarizes the hardware specifications of the host machine and the container execution runtime.

\begin{table}[H]
\centering
\small
\caption{Experimental Host Hardware and Container Execution Environment}
\label{tab:test-hardware-environment}
\vspace{0.2cm}
\begin{tabular}{ll}
\toprule
\textbf{System Parameter} & \textbf{Specification} \\
\midrule
Host Processor (CPU)      & Intel(R) Core(TM) Ultra 7 265H (16 physical cores, x86\_64) \\
CPU Cache Hierarchy       & 24~MiB L3 cache, 48~MiB L2 cache, 1.75~MiB L1 cache \\
Physical Host Memory      & 64~GB physical RAM \\
Dedicated Container RAM   & 31~GiB allocated memory with 8.0~GiB swap (WSL2 managed) \\
Local Storage             & 1~TB NVMe SSD (\texttt{ext4} virtual disk) \\
Host Operating System     & Microsoft Windows 11 Enterprise / Professional \\
Linux Runtime Subsystem   & Ubuntu 24.04.4 LTS (Noble Numbat) on WSL2 (Kernel 6.6.87) \\
Container Engine          & Docker Engine 29.5.3-rd \\
Orchestration Tool        & Docker Compose v5.1.4 \\
\bottomrule
\end{tabular}
\end{table}

The virtualization architecture utilizes the Windows Subsystem for Linux (WSL2) backed by a lightweight Hyper-V hypervisor. Unlike traditional full-system virtual machines, WSL2 provides direct native hardware access to all 16 physical CPU cores without core pinning overhead, while dynamically allocating a bounded 31~GiB RAM slice dedicated to Docker containers. This ensures that the compute and memory benchmarks faithfully represent real-world bare-metal industrial edge deployment conditions.

\clearpage

\section{Testing: testing adjusted simulator}
\label{sec:testing}

% !TeX root = ../main.tex
% Placeholder reused by every check of this section. Drop the screenshot at
% figures/simValidationCaptureRun.png and all the figures of the section update at
% once; while the file is missing, the example-image shipped with the mwe package is
% typeset instead so the document keeps compiling.
\newcommand{\physfig}[2]{%
  \par\vspace{6pt}%
  \begin{center}%
    \IfFileExists{figures/#1.png}{%
      \includegraphics[width=0.85\linewidth]{#1}%
    }{%
      \includegraphics[width=0.85\linewidth]{example-image}%
    }\\[2pt]%
    \captionof{figure}{#2}\label{fig:#1}%
  \end{center}%
  \vspace{2pt}%
}

The physical powertrain model formulated in Section~\ref{sec:physical_modeling} and incorporated into the simulator as described in Section~\ref{sec:synthetic_dataset_generation} replaces heuristic signal
generation with an explicit electromechanical model: an armature circuit, a torque balance
referred to the motor shaft, a lumped thermal node for the winding, a battery with an
equivalent-circuit terminal voltage and a set of transducers that inject noise,
quantisation and drift into every published channel. Introducing that model raises a
question that unit tests alone cannot answer: whether the telemetry that finally
reaches the historian still \emph{looks} like the physical system it claims to
represent.

\textbf{Nature of this test.} The checks collected in this section are
\emph{functional, informational tests} executed by \textbf{manual inspection} of the
telemetry rendered in the Grafana dashboards. They are not automated tests, they are
not executed by a CI/CD pipeline, they produce no pass/fail artefact and they are not
regression protected: every run is observed, interpreted and archived by hand. Their
purpose is evidential rather than verificational, to show a reader of the thesis that
the simulated signals reproduce the qualitative and quantitative behaviour expected
from the electromechanical model, so that the anomalies synthesised later on can be
argued to be physically plausible and not artefacts of arbitrary signal shaping.

\subsection{Test Environment and Procedure}

All observations were taken from the deployed stack: the simulated fleet publishes one
VDA~5050 \texttt{state} message per robot and second, which travels through the MQTT broker
through Kafka Connect into the QuestDB table \texttt{vda5050\_state}, where
the Grafana dashboards, representations of each source form read them. The state publication rate of one hertz is
the effective sampling rate of every observation reported below and sets the limit of
what can be seen.
A gaussian noise is introduce to simulate the real world behaviour.

\subsection{Motor voltage and governing variables (Vm, Im, wm, Ra, La, Ke)}
To verify that the simulator adheres to the electromechanical motor model, the quasi-steady armature equation $V_m(t) \approx R_a(T_m) \cdot I_m(t) + K_e \cdot \omega_m(t)$ (neglecting the fast inductive transient $L_a \frac{dI_m}{dt} \approx 0$) was evaluated using telemetry sampled during cruise motion. With an armature resistance $R_a = 0.0501\,\Omega$, current $I_m = 2.9688\,\text{A}$, rotor speed $\omega_m = 100.0413\,\text{rad/s}$, and $K_e = 0.19\,\text{V}\cdot\text{s/rad}$, the calculated back-EMF is $E_b = K_e \cdot \omega_m = 19.0078\,\text{V}$ (closely matching the measured $E_{b,\text{meas}} = 18.9922\,\text{V}$) and the ohmic drop is $R_a \cdot I_m = 0.1487\,\text{V}$. The theoretical terminal voltage yields $V_{m,\text{calc}} = 0.1487\,\text{V} + 18.9922\,\text{V} = 19.1409\,\text{V}$, which deviates by only $0.0550\,\text{V}$ (a relative error of $0.28\%$) from the measured $V_{m,\text{meas}} = 19.0859\,\text{V}$. This minor discrepancy is completely bounded by the simulated 12-bit ADC quantization step ($\approx 7.8\,\text{mV}$) and the Gaussian sensor noise ($\sigma_v = 0.03\,\text{V}$), confirming that the simulator accurately follows the physical motor principle.
\physfig{motor_voltage_form}{Motor voltage form, with all the terms}

\subsection{Motor Torque and governing variables (\texorpdfstring{$\tau_m$, $K_t$, $I_m$}{tau\_m, Kt, Im})}
To verify that the simulator accurately governs electromagnetic torque generation, the linear torque-current relationship $\tau_m(t) = K_t \cdot I_m(t)$ of Equation~\eqref{eq:electromagnetic_torque} was evaluated against the Grafana telemetry panel shown in Figure~\ref{fig:torque_form}. At the sampled operating point ($t = \text{15:56:38}$, during an idle/standby interval before motion), the torque constant is $K_t = 0.1900\,\text{N}\cdot\text{m/A}$ (numerically equivalent to the back-EMF constant $K_e = 0.1900\,\text{V}\cdot\text{s/rad}$) and the armature current is $I_m = 0.0391\,\text{A}$. The calculated electromagnetic torque yields $\tau_{m,\text{calc}} = K_t \cdot I_m = 0.1900\,\text{N}\cdot\text{m/A} \times 0.0391\,\text{A} = 0.007429\,\text{N}\cdot\text{m}$, matching the dashboard's calculated channel (\texttt{calculated\_torque\_nm} = $0.0074\,\text{N}\cdot\text{m}$) to four decimal places. The measured motor torque reported by the drive telemetry is $\tau_{m,\text{meas}} = 0.0077\,\text{N}\cdot\text{m}$ (\texttt{motor\_torque\_nm}), resulting in an absolute discrepancy of only $0.0003\,\text{N}\cdot\text{m}$ ($0.3\,\text{mN}\cdot\text{m}$). This minute deviation is completely bounded by the 12-bit ADC quantization resolution and the Gaussian transducer noise ($\sigma_I \approx 0.08\,\text{A}$, translating to $\sigma_\tau \approx 0.015\,\text{N}\cdot\text{m}$). Across the entire mission profile, dynamic tracking remains consistent: during acceleration ramps, current surges to $I_m \approx 8.63\,\text{A}$ producing a torque peak of $\approx 1.63\,\text{N}\cdot\text{m}$, whereas during steady-state cruise at $I_m \approx 3.0$--$3.2\,\text{A}$, torque settles stably at $\approx 0.57$--$0.61\,\text{N}\cdot\text{m}$. The measured and calculated torque traces overlap seamlessly throughout all motion phases, confirming that the simulated machine faithfully reproduces electromagnetic torque conversion without synthetic distortions.
\physfig{torque_form}{Motor torque form, with all the terms}

\subsection{Mechanical torque balance and load dynamics (\texorpdfstring{$\tau_m$, $\tau_L$, $\omega_m$, $B_m$, $J_m$}{tau\_m, tau\_L, wm, Bm, Jm})}
The dynamic torque equilibrium referred to the motor rotor shaft is governed by Equation~\eqref{eq:torque_balance}, $J_m \frac{d\omega_m(t)}{dt} = \tau_m(t) - B_m \cdot \omega_m(t) - \tau_{L,\text{ref}}(t)$, or expressed in terms of the total mechanical demand: $\tau_{\text{demand}}(t) = J_m \frac{d\omega_m(t)}{dt} + B_m \cdot \omega_m(t) + \tau_{L,\text{ref}}(t)$, where $\tau_{\text{viscous}}(t) = B_m \cdot \omega_m(t)$ is the internal viscous friction torque and $\tau_{L,\text{ref}}(t) = \frac{\tau_L(t)}{\eta_g i_g}$ represents the ground load torque reflected through the reduction gearbox ($i_g = 10$, $\eta_g = 0.85$). This formulation was verified using the Grafana telemetry panel in Figure~\ref{fig:mechanicalTorqueBalance} sampled during steady cruise motion ($t = \text{15:57:16}$). At this operating point, the rotor speed is stabilized at $\omega_m = 99.9031\,\text{rad/s}$ (corresponding to the vehicle travelling at cruise velocity $v \approx 1.0\,\text{m/s}$), rendering the angular acceleration null ($\frac{d\omega_m}{dt} \approx 0$) and nullifying the inertial demand ($J_m \dot{\omega}_m \approx 0$). The reported viscous friction torque is $\tau_{\text{viscous}} = 0.1499\,\text{N}\cdot\text{m}$, which confirms a viscous damping coefficient $B_m = \frac{\tau_{\text{viscous}}}{\omega_m} = \frac{0.1499}{99.9031} \approx 0.0015\,\text{N}\cdot\text{m}\cdot\text{s/rad}$. The reflected load torque is $\tau_{L,\text{ref}} = 0.4155\,\text{N}\cdot\text{m}$, corresponding to the rolling resistance of the $180\,\text{kg}$ vehicle on level ground ($F_{\text{roll}} = C_{rr} m_{\text{total}} g = 0.02 \times 180\,\text{kg} \times 9.81\,\text{m/s}^2 \approx 35.32\,\text{N}$, giving $\tau_L = F_{\text{roll}} \cdot r_w \approx 3.532\,\text{N}\cdot\text{m}$ at wheel radius $r_w = 0.10\,\text{m}$, and $\tau_{L,\text{ref}} = \frac{3.532}{0.85 \times 10} \approx 0.4155\,\text{N}\cdot\text{m}$). The theoretical total mechanical demand thus evaluates to $\tau_{\text{demand}} = \tau_{L,\text{ref}} + \tau_{\text{viscous}} = 0.4155\,\text{N}\cdot\text{m} + 0.1499\,\text{N}\cdot\text{m} = 0.5654\,\text{N}\cdot\text{m}$, perfectly coinciding with the reported channel \texttt{total\_mechanical\_demand\_nm} ($0.5654\,\text{N}\cdot\text{m}$). The electromagnetic motor torque delivered by the drive at this moment is $\tau_{m,\text{meas}} = 0.6130\,\text{N}\cdot\text{m}$, presenting a minor difference of $0.0476\,\text{N}\cdot\text{m}$ attributable to simulated current transducer noise ($\sigma_\tau \approx 0.015\,\text{N}\cdot\text{m}$) and speed controller regulation ripple. During acceleration ramps, the inertial term $J_m\dot{\omega}_m$ momentarily elevates the torque demand to $\approx 1.63\,\text{N}\cdot\text{m}$, validating that the simulator accurately integrates both static resistive loads and dynamic inertial terms.
\physfig{mechanicalTorqueBalance}{Mechanical torque balance and load dynamics, with all the terms}

\subsection{Kinematic closure and wheel slip dynamics (\texorpdfstring{$v$, $\omega_m$, $s$, $r_w$, $i_g$}{v, wm, s, rw, ig})}
The kinematic coupling between the electric motor shaft, transmission, driven wheels, and vehicle ground translation is defined by the kinematic closure relation $v(t) = \frac{\omega_m(t)}{i_g} \cdot r_w \cdot (1 - s(t))$, where $v_w(t) = \frac{\omega_m(t)}{i_g} \cdot r_w$ represents the theoretical peripheral linear velocity of the wheel, $i_g = 10$ is the transmission reduction ratio, $r_w = 0.10\,\text{m}$ is the nominal wheel radius, and $s(t)$ is the longitudinal wheel slip ratio. This formulation was validated using the Grafana telemetry panel shown in Figure~\ref{fig:linearVelocityForm} sampled during steady cruise motion ($t = \text{15:56:57}$). At this operating point, the reported wheel linear speed is $v_w = 1.0011\,\text{m/s}$ (corresponding to an angular motor velocity $\omega_m = \frac{v_w i_g}{r_w} = 100.11\,\text{rad/s}$), and the estimated longitudinal wheel slip is $s = -0.0073$ (a slight negative micro-slip of $-0.73\%$ induced by transducer Gaussian noise and nominal tire contact compliance). Evaluating the kinematic closure equation yields a theoretical ground speed of $v_{\text{calc}} = v_w \cdot (1 - s) = 1.0011\,\text{m/s} \cdot (1 - (-0.0073)) = 1.0084\,\text{m/s}$, matching the dashboard's calculated channel (\texttt{calculated\_ground\_speed\_v} = $1.0084\,\text{m/s}$) to four decimal places. When compared to the measured ground speed $v_{\text{meas}} = 1.0000\,\text{m/s}$ (\texttt{ground\_speed\_v}), the absolute discrepancy is only $0.0084\,\text{m/s}$ (a relative error of $0.84\%$). Across the mission profile, the wheel speed tracks the ground translation with minimal slip ($s \approx 0$) during unperturbed cruise, while the brief velocity reductions at navigation waypoints are tracked synchronously by both channels. This confirms that the transmission kinematics and slip interaction are physically consistent and free of numerical drift.
\physfig{linearVelocityForm}{Linear velocity form, with all the terms}

\subsection{Motor thermal dynamics and power loss balance (\texorpdfstring{$T_m$, $T_{\mathrm{amb}}$, $P_{\mathrm{joule}}$, $P_{\mathrm{iron}}$, $P_{\mathrm{cool}}$}{Tm, Tamb, Pjoule, Piron, Pcool})}
The thermal evolution of the motor winding is governed by the lumped first-order power balance $C_{th} \frac{dT_m(t)}{dt} = P_{\text{joule}}(t) + P_{\text{iron}}(t) - P_{\text{cool}}(t)$, where $P_{\text{joule}}(t) = I_m(t)^2 R_a(T_m)$ represents the resistive copper losses, $P_{\text{iron}}(t) = C_{fe} \omega_m(t)^{1.6} B_{\max}^2$ represents core iron losses, and $P_{\text{cool}}(t) = \frac{T_m(t) - T_{\text{amb}}}{R_{th}}$ is the convective heat transfer to the ambient environment ($T_{\text{amb}} = 22.0000\,^{\circ}\text{C}$). This thermal balance was evaluated using the telemetry panel shown in Figure~\ref{fig:motor_termical_balance} at $t = \text{15:57:07}$ during an acceleration transition between navigation nodes. At this instant, the telemetry records copper losses of $P_{\text{joule}} = 3.6831\,\text{W}$ and core iron losses of $P_{\text{iron}} = 3.7758\,\text{W}$, yielding a total heat dissipation rate of $P_{\text{total}} = 3.6831\,\text{W} + 3.7758\,\text{W} = 7.4589\,\text{W}$. With the winding temperature at $T_m = 23.1900\,^{\circ}\text{C}$, the temperature rise above ambient is $T_m - T_{\text{amb}} = 1.1900\,\text{K}$, producing a cooling power of $P_{\text{cool}} = 0.4760\,\text{W}$, which exactly identifies the thermal resistance $R_{th} = \frac{T_m - T_{\text{amb}}}{P_{\text{cool}}} = \frac{1.1900\,\text{K}}{0.4760\,\text{W}} = 2.5000\,\text{K/W}$. The instantaneous asymptotic equilibrium temperature is given by $T_{\infty} = T_{\text{amb}} + P_{\text{total}} \cdot R_{th} = 22.0000\,^{\circ}\text{C} + (7.4589\,\text{W} \times 2.5000\,\text{K/W}) = 40.64725\,^{\circ}\text{C} \approx 40.6473\,^{\circ}\text{C}$, which matches the dashboard channel \texttt{equilibrium\_temp\_c} ($40.6473\,^{\circ}\text{C}$) to four decimal places. The net thermal power $P_{\text{net}} = P_{\text{total}} - P_{\text{cool}} = 6.9829\,\text{W}$ is integrated across the thermal capacitance $C_{th} = 80\,\text{J/K}$, producing a smooth, continuous exponential temperature rise ($\frac{dT_m}{dt} \approx 0.087\,\text{K/s}$) without artificial discontinuities, verifying the physical validity of the thermal node.
\physfig{motor_termical_balance}{Motor thermal dynamic and power loss balance, with all the terms}

\subsection{Battery terminal voltage and internal-resistance drop (\texorpdfstring{$V_{\mathrm{bat}}$, $V_{\mathrm{ocv}}$, $R_{\mathrm{int}}$, $I_{\mathrm{bat}}$}{Vbat, Vocv, Rint, Ibat})}
The terminal voltage of the traction battery is governed by the simplified equivalent-circuit model (ECM) of Equation~\eqref{eq:battery_ecm}: $V_{\text{bat}}(t) = V_{\text{ocv}}(\text{SOC}) - \Delta V_{\text{ohmic}}(t)$, where $\Delta V_{\text{ohmic}}(t) = I_{\text{bat}}(t) \cdot R_{\text{int}}(\text{SOC}, T_{\text{bat}})$ represents the instantaneous ohmic drop across the pack internal resistance. This relation was examined using the Grafana telemetry panel shown in Figure~\ref{fig:battery_tension_form} sampled during cruise motion ($t = \text{15:57:11}$). At this point, the pack delivers a current of $I_{\text{bat}} = 3.0273\,\text{A}$ and exhibits an internal resistance of $R_{\text{int}} = 0.0535\,\Omega$. The theoretical ohmic sag evaluates to $\Delta V_{\text{ohmic}} = I_{\text{bat}} \cdot R_{\text{int}} = 3.0273\,\text{A} \times 0.0535\,\Omega = 0.16196\,\text{V} \approx 0.1619\,\text{V}$, exactly matching the dashboard's reported channel (\texttt{ohmic\_sag\_v} = $0.1619\,\text{V}$). Subtracting this drop from the open-circuit voltage $V_{\text{ocv}} = 23.2248\,\text{V}$ yields a theoretical terminal voltage of $V_{\text{bat},\text{calc}} = 23.2248\,\text{V} - 0.16196\,\text{V} = 23.06284\,\text{V} \approx 23.0629\,\text{V}$, coinciding with \texttt{v\_bat\_calc} ($23.0629\,\text{V}$) to four decimal places. The measured terminal voltage reported by telemetry is $V_{\text{bat},\text{meas}} = 23.0938\,\text{V}$ (\texttt{v\_bat\_meas}), differing by only $0.0309\,\text{V}$ (a relative error of $0.13\%$). This minor discrepancy is completely bounded by the simulated 12-bit ADC quantization step ($\approx 7.8\,\text{mV}$) and the battery voltage sensor Gaussian noise ($\sigma_v = 0.03\,\text{V}$). Dynamically, the panel demonstrates the expected electrochemical response: while idling ($I_{\text{bat}} \approx 0.5\,\text{A}$), the ohmic drop remains minimal and $V_{\text{bat}} \approx V_{\text{ocv}}$; under cruise load ($I_{\text{bat}} \approx 3.0\,\text{A}$), terminal voltage experiences a sustained sag of $\approx 0.16\,\text{V}$; and during acceleration surges ($I_{\text{bat}} \approx 6.8\,\text{A}$), the ohmic dip deepens to $\approx 0.36\,\text{V}$ before recovering, fully validating the ECM implementation.
\physfig{battery_tension_form}{Battery tension form, with all terms}

\subsection{Interpretation and Limitations}

The values quoted above are consistency checks against the model, not metrological
measurements: every one of them is conditioned by the calibration constants published in
the simulator configuration, and they hold for the operating point described, a
one-hundred-and-eighty-kilogram vehicle travelling unloaded on level ground at
$1.0~\mathrm{m/s}$. Changing the payload, the grade or the speed moves the numerical
values while leaving the relations intact, which is precisely the property the checks
exercise.

The one hertz publication interval also bounds the scope of the section. Electromagnetic
transients of tens of milliseconds, the ripple harmonics injected by the stator
short-circuit fault, and the slow drift of the thermistor and the current transducer,
specified in fractions of a degree or a percent per hour, are outside what can be
observed here; they are asserted deterministically in the automated test suite of the
simulator instead. This section is consequently an informational, human-reviewed
characterisation of the simulated telemetry, whose role in this work is to justify that
the data on which the anomaly detection models are trained and evaluated behaves like a
powertrain.


\section{Evaluation: Data pipeline evaluation}
\label{sec:data_pipeline_evaluation}

% !TeX root = ../main.tex
\providecolor{primary}{RGB}{24, 76, 120}
\providecolor{secondary}{RGB}{198, 40, 40}
\providecolor{accentgreen}{RGB}{46, 125, 50}
\providecolor{darkgray}{RGB}{60, 60, 60}

This section assesses the operational resilience, throughput stability, and resource consumption of the end-to-end industrial telemetry streaming pipeline under continuous real-time load. In distributed IoT streaming architectures, standard unit and integration tests are insufficient to guarantee production readiness; long-running \textit{soak tests} are indispensable to expose insidious failure modes such as progressive memory leaks, queue backpressure, partition fragmentation, and metadata accumulation.

%======================================================================
\subsection{Evaluation Methodology and Workload Design}
\label{subsec:eval_methodology}

The evaluation campaign was designed to monitor the complete containerized platform over an extended multi-day lifecycle. Across the entire experimental period, the cluster operated without interruption for \textbf{62.35~hours (2.6~days)}---spanning from the first message published on 2026-11-09 at 17:29:29 UTC to 2026-11-12 at 07:50:17 UTC.

\paragraph{Synthetic Industrial Workload.}
The streaming workload was continuously produced by a simulated fleet of 5 mobile industrial robots (AGVs/AMRs) operating with the electromechanical vehicle simulator described in Section~\ref{sec:source_system_simulation}. Each simulated robot published state messages over MQTT at a constant frequency of 1~Hz (totaling a continuous platform ingestion rate of 5~messages/second). Messages traveled through the Eclipse Mosquitto broker, were bridged by Telegraf into the Redpanda \texttt{vda5050-state} topic, transformed in real time by Apache Spark Structured Streaming, persisted into the Apache Iceberg lakehouse on MinIO S3, and ingested into QuestDB for real-time dashboard visualization.

\paragraph{Unattended Overnight Soak Monitoring Protocol.}
To eliminate artificial interference caused by interactive desktop usage and operator queries, the empirical auditing sessions were conducted during unattended overnight windows when the workstation was left dedicated entirely to continuous ingestion:
\begin{enumerate}
    \item \textbf{Phase~1 (Day~1 Baseline Soak, 14.30~Hours):} An initial overnight audit (from 2026-11-09 18:31:02 to 2026-11-10 08:48:49 UTC, $N = 1,641$ snapshots) executed under default streaming settings. This phase established the baseline performance and successfully uncovered critical operational vulnerabilities.
    \item \textbf{Phase~2 (Days~2--3 Extended Validation Soak, 38.86~Hours):} Following daytime root-cause analysis and hardening deployment, an extended multi-day overnight soak test was launched (from 2026-11-10 17:59:59 to 2026-11-12 08:49:44 UTC, $N = 4,450$ snapshots). This phase monitored the hardened platform continuously across two consecutive nights and the intervening day to confirm asymptotic steady-state stability.
\end{enumerate}

Throughout both sessions, an automated background daemon captured discrete \texttt{docker stats} snapshots every 30~seconds across all 11 containerized services, recording CPU utilization, resident set size (RSS) memory, network I/O, block I/O, and thread counts.

%======================================================================
\subsection{Phase 1: Baseline Soak Test and Bottleneck Discovery (Day 1)}
\label{subsec:eval_phase1}

\paragraph{Objective.}
The primary objective of the initial baseline soak test was to verify whether the streaming pipeline could sustain continuous processing without worker crashes, pipeline freezes, or message lag accumulation under default Apache Iceberg configurations.

\paragraph{Baseline Empirical Findings.}
During the 14.30-hour initial run, the platform achieved 100\% service availability without crashes, restarts, or unhandled exceptions. Ingested data integrity was verified at the storage sinks: exactly 233,412 rows were persisted into QuestDB (\texttt{fleet\_features}) with uninterrupted timestamps. Furthermore, Apache Spark maintained strict real-time parity with Redpanda without backpressure:
\begin{equation}
    \text{avgOffsetsBehindLatest} = 0.0, \quad \text{maxOffsetsBehindLatest} = 0
\end{equation}
These streaming metrics quantify consumer lag by tracking the offset difference between the latest message published to the broker and the offset committed by Spark. A zero value across the entire soak test proves that micro-batches were processed instantaneously upon arrival without latency drift.

Table~\ref{tab:pipeline_perf_metrics_phase1} details the container resource utilization recorded across the 14.30-hour baseline execution.

\begin{table}[htbp]
\centering
\scriptsize
\caption{Container Resource Utilization and Memory Drift (Phase 1: Day 1 Baseline Soak Test)}
\label{tab:pipeline_perf_metrics_phase1}
\vspace{0.2cm}
\resizebox{\textwidth}{!}{
\begin{tabular}{l r r r r r r r l}
\toprule
\textbf{Service / Container} & \textbf{CPU$_{\text{avg}}$} & \textbf{CPU$_{\text{max}}$} & \textbf{MEM$_{\text{avg}}$} & \textbf{MEM$_{\text{max}}$} & \textbf{MEM$_{\text{Start}}$} & \textbf{MEM$_{\text{End}}$} & \textbf{$\Delta$ MEM} & \textbf{Observed Trend} \\
\midrule
\texttt{spark-fleet-features}     & 24.3\% & 131.9\% & 2.1 GiB & 3.1 GiB & 2.9 GiB & 1.7 GiB & -1.2 GiB & Stabilized (GC) \\
\texttt{spark-vda5050-state}      &  2.5\% &  48.5\% & 1.6 GiB & 2.1 GiB & 2.1 GiB & 1.3 GiB & -767 MiB & Stabilized (GC) \\
\texttt{questdb}                  & 12.8\% &  27.2\% & 569 MiB & 741 MiB & 711 MiB & 425 MiB & -287 MiB & Decreased (Stable) \\
\texttt{redpanda}                 &  1.3\% &   9.1\% & 1.8 GiB & 2.1 GiB & 1.9 GiB & 1.4 GiB & -515 MiB & Stabilized \\
\texttt{minio}                    &  5.8\% &  33.2\% & 1.2 GiB & 2.1 GiB & 400 MiB & 2.0 GiB & \textbf{+1.6 GiB} & \textbf{+406\% (Metadata Cache)} \\
\texttt{questdb-kafka-connect}    &  0.9\% &  15.2\% & 581 MiB & 939 MiB & 616 MiB & 476 MiB & -141 MiB & Stabilized \\
\texttt{grafana}                  &  0.6\% &   6.3\% & 150 MiB & 173 MiB & 150 MiB & 148 MiB &   -1 MiB & Flat \\
\texttt{simulator} (Rust AGV)     &  0.5\% &   1.1\% & 6.1 MiB & 6.7 MiB & 4.2 MiB & 6.3 MiB &   +2 MiB & Negligible (Flat) \\
\texttt{telegraf}                 &  0.2\% &   1.5\% &  60 MiB & 107 MiB &  72 MiB &  44 MiB &  -28 MiB & Stabilized \\
\texttt{mosquitto}                &  0.5\% &   5.9\% & 6.9 MiB &  12 MiB & 7.9 MiB & 6.4 MiB &   -1 MiB & Flat \\
\texttt{sim-control-panel}        &  0.1\% &   0.2\% &  27 MiB &  28 MiB &  25 MiB &  26 MiB &   +1 MiB & Flat \\
\bottomrule
\end{tabular}
}
\end{table}

\paragraph{Critical Bottleneck Discovery: Unbounded Metadata Accumulation.}
While computational streaming components behaved stably, post-test inspection revealed severe operational anomalies centered on the storage tier:
\begin{enumerate}
    \item \textbf{Monotonic MinIO RAM Growth (+406\%):} As illustrated in Figure~\ref{fig:pipeline_memory_drift_phase1}, MinIO resident memory diverged sharply from all other services, climbing continuously from 400~MiB up to 2.07~GiB (+1.6~GiB).
    \item \textbf{Object Storage Exhaustion (128~GiB):} The MinIO bucket reached 128~GiB of disk usage. Crucially, actual persisted Parquet telemetry data totaled only 321~MB; the remaining \textbf{127.5~GiB (99.7\% of the entire bucket)} consisted of over 61,000 Apache Iceberg metadata JSON files.
\end{enumerate}

\begin{figure}[htbp]
\centering
\begin{tikzpicture}
\begin{axis}[
    title={\textbf{Phase 1: Memory Footprint Over 14.3 Hours (Baseline Execution)}},
    xlabel={Execution Time (Hours)},
    ylabel={RAM Consumption (GiB)},
    xmin=0, xmax=14,
    ymin=0, ymax=3.5,
    width=0.92\textwidth,
    height=6.2cm,
    grid=both,
    grid style={line width=.1pt, draw=gray!20},
    legend style={at={(0.98,0.95)}, anchor=north east, font=\scriptsize},
    thick
]
    % Spark Fleet Features
    \addplot[color=primary, mark=square*] coordinates {
        (0, 2.31) (2, 2.56) (4, 2.22) (6, 2.07) (8, 1.93) (10, 1.79) (12, 1.70) (14, 1.65)
    };
    \addlegendentry{Spark Fleet Features (Stateless JVM)}

    % MinIO (Object Storage)
    \addplot[color=secondary, mark=triangle*] coordinates {
        (0, 0.19) (2, 0.59) (4, 0.86) (6, 1.10) (8, 1.37) (10, 1.53) (12, 1.78) (14, 2.07)
    };
    \addlegendentry{MinIO S3 (Metadata Index Drift)}

    % Redpanda
    \addplot[color=accentgreen, mark=diamond*] coordinates {
        (0, 1.90) (2, 1.97) (4, 2.02) (6, 2.02) (8, 1.86) (10, 1.65) (12, 1.48) (14, 1.41)
    };
    \addlegendentry{Redpanda Kafka (Seastar Thread-Pool)}

    % QuestDB
    \addplot[color=orange, mark=circle*] coordinates {
        (0, 0.69) (2, 0.60) (4, 0.60) (6, 0.60) (8, 0.56) (10, 0.50) (12, 0.44) (14, 0.42)
    };
    \addlegendentry{QuestDB (Zero-GC Columnar Engine)}
\end{axis}
\end{tikzpicture}
\caption{Evolution of memory footprint during the Phase~1 baseline soak test. While Spark, Redpanda, and QuestDB stabilized downward, MinIO exhibited runaway monotonic growth due to accumulating Iceberg metadata indices.}
\label{fig:pipeline_memory_drift_phase1}
\end{figure}

%======================================================================
\subsection{Root-Cause Analysis and Lakehouse Pipeline Hardening}
\label{subsec:eval_hardening}

\paragraph{Mechanism of the Failure Mode.}
In Apache Iceberg, every streaming micro-batch commit generates a new metadata version file ($v_N\text{.metadata.json}$). By default, the specification appends the entire historical snapshot log into every newly generated metadata file without purging previous versions. Consequently, storage consumption scales quadratically with the cumulative number of commits, causing individual metadata files to expand progressively throughout the run:
\begin{itemize}
    \item Early version ($v_1$): 12~KB.
    \item Intermediate version ($v_{1000}$): 872~KB.
    \item Mid-test version ($v_{10000}$): 8.3~MB.
    \item End-of-night version ($v_{17650}$): 15.0~MB.
\end{itemize}
With micro-batches executing every 3~seconds, Spark was serializing and transmitting a 15~MB JSON file to MinIO every 3~seconds. This inflated Spark commit latency from 180~ms to 3--5~seconds, saturated network I/O, and forced the MinIO daemon to retain massive inode structures and bloom filters in RAM, explaining the observed 2.07~GiB memory drift.

Furthermore, executing standard maintenance via \texttt{CALL expire\_snapshots} only deleted orphaned Parquet data files; table-level metadata JSON files were permanently preserved as audit artifacts under default table properties.

\paragraph{Remediation Strategy: Table Property Hardening and Smart Watcher.}
To eliminate this vulnerability permanently, a two-pronged architectural remediation was implemented:
\begin{enumerate}
    \item \textbf{Iceberg Table Schema Properties:} The DDL definitions in \texttt{src/schemas/fleet\_features.py} and \texttt{src/schemas/vda5050\_state.py} were injected with automatic commit-time pruning:
    \begin{itemize}
        \item \texttt{'write.metadata.delete-after-commit.enabled' = 'true'}: Instructs Spark to automatically purge obsolete metadata files synchronously upon committing new snapshots.
        \item \texttt{'write.metadata.previous-versions-max' = '5'}: Restricts the retained metadata history to the 5 most recent version files, bounding metadata file size permanently below 40~KB.
    \end{itemize}
    \item \textbf{Decoupled Smart Maintenance Watcher Daemon:} An autonomous background service (\texttt{maintenance\_watcher.py}) was deployed to monitor MinIO S3 metadata via lightweight HTTP requests (without instantiating a heavy Spark JVM). Whenever accumulated snapshots breach 600 or metadata size exceeds 500~KB, the daemon executes an isolated \texttt{spark-submit} maintenance job to compact Parquet files and expire snapshots, fully releasing its memory upon task completion.
\end{enumerate}

Following an initial administrative purge of historical orphan metadata, the immediate impact of the hardening configuration is summarized in Table~\ref{tab:pipeline_remediation}.

\begin{table}[htbp]
\centering
\small
\caption{Storage and Pipeline Performance Metrics Before vs. After Metadata Hardening}
\label{tab:pipeline_remediation}
\vspace{0.2cm}
\begin{tabular}{l r r r}
\toprule
\textbf{Metric} & \textbf{Pre-Hardening (Day 1)} & \textbf{Post-Hardening} & \textbf{Improvement} \\
\midrule
Total Lakehouse S3 Storage & 128.0 GiB & \textbf{36.0 MiB} & \textbf{-99.97\%} \\
Silver Table Metadata Size & 122.5 GiB & \textbf{26.0 MiB} & -99.98\% \\
Bronze Table Metadata Size &   5.0 GiB &  \textbf{1.7 MiB} & -99.96\% \\
Per-Commit Metadata JSON   &  15.0 MiB &  \textbf{8.7 KiB} & -99.94\% \\
Retained Metadata Files    & 17,654 versions & \textbf{5 versions} & Bound \& Static \\
Spark Iceberg Commit Time  & 2.5s -- 5.0s & \textbf{0.18s} & \textbf{15$\times$ Faster} \\
\bottomrule
\end{tabular}
\end{table}

%======================================================================
\subsection{Phase 2: Extended Multi-Day Validation Soak Test (Days 2--3)}
\label{subsec:eval_phase2}

\paragraph{Objective.}
The objective of Phase~2 was to conduct an extended multi-day soak test spanning across two consecutive nights (Days~2 and 3) to empirically validate that the deployed remediations permanently bound metadata expansion, eliminate resource drift, and maintain sub-second streaming latencies under extended continuous operation.

\paragraph{Multi-Day Operational Findings.}
The Phase~2 soak test executed autonomously for \textbf{38.86~hours} (from 2026-11-10 17:59:59 to 2026-11-12 08:49:44 UTC), logging $N = 4,450$ discrete measurement snapshots. Combined with the baseline run and daytime operation, the cluster achieved \textbf{62.35~hours of uninterrupted streaming}:
\begin{itemize}
    \item \textbf{Milestone Ingestion Scale:} Total records persisted into QuestDB exceeded \textbf{1,026,060 rows} (> 1 million events) spanning from 2026-11-09 to 2026-11-12 with zero dropped messages and zero streaming lag ($\text{avgOffsetsBehindLatest} = 0.0$).
    \item \textbf{Autonomous Compaction Performance:} The background maintenance watcher daemon triggered exactly \textbf{87 consecutive optimization cycles} across the 39-hour run (activating punctually every $\sim30$~minutes at 600 snapshots). Each cycle executed in an average of 26.2~seconds with exit code 0.
    \item \textbf{Bounded Storage Verification:} Total lakehouse storage on MinIO S3 concluded at only \textbf{70 MiB} across 1,073 objects (compared to 128 GiB and 61,000 files in Phase~1), definitively proving that storage consumption is permanently bounded.
\end{itemize}

Table~\ref{tab:pipeline_perf_metrics_phase2} details container resource utilization across the extended 38.86-hour execution.

\begin{table}[htbp]
\centering
\scriptsize
\caption{Container Resource Utilization and Memory Drift (Phase 2: Days 2--3 Extended 38.9-Hour Soak Test)}
\label{tab:pipeline_perf_metrics_phase2}
\vspace{0.2cm}
\resizebox{\textwidth}{!}{
\begin{tabular}{l r r r r r r r l}
\toprule
\textbf{Service / Container} & \textbf{CPU$_{\text{avg}}$} & \textbf{CPU$_{\text{max}}$} & \textbf{MEM$_{\text{avg}}$} & \textbf{MEM$_{\text{max}}$} & \textbf{MEM$_{\text{Start}}$} & \textbf{MEM$_{\text{End}}$} & \textbf{$\Delta$ MEM} & \textbf{Observed Trend} \\
\midrule
\texttt{spark-fleet-features}     & 14.3\% & 109.6\% & 2.2 GiB & 2.4 GiB & 2.2 GiB & 2.1 GiB & -160 MiB & Stabilized (-7\% GC) \\
\texttt{spark-vda5050-state}      &  2.4\% &  58.0\% & 1.5 GiB & 1.5 GiB & 1.4 GiB & 1.5 GiB & +140 MiB & Stabilized (GC) \\
\texttt{questdb}                  & 17.5\% & 138.5\% & 1.5 GiB & 1.6 GiB & 1.5 GiB & 1.6 GiB &  +88 MiB & Stable (mmap cache) \\
\texttt{redpanda}                 &  1.2\% &   9.2\% & 2.4 GiB & 3.1 GiB & 1.7 GiB & 3.1 GiB & +1.4 GiB & Bounded (Seastar Pool) \\
\texttt{minio}                    &  4.1\% & 266.3\% & 690 MiB & 808 MiB & 512 MiB & 782 MiB & +270 MiB & Plateau (Kernel Slab) \\
\texttt{spark-maintenance}        &  4.8\% & 898.1\% & 382 MiB & 2.6 GiB & 380 MiB & 405 MiB &  +24 MiB & Stable (Periodic Daemon) \\
\texttt{questdb-kafka-connect}    &  0.9\% &   8.2\% & 883 MiB & 991 MiB & 862 MiB & 857 MiB &   -6 MiB & Stable (JVM) \\
\texttt{grafana}                  &  0.6\% &  14.9\% & 197 MiB & 833 MiB & 171 MiB & 201 MiB &  +30 MiB & Stable (Sampled Queries) \\
\texttt{simulator} (Rust AGV)     &  0.6\% &   2.5\% & 6.3 MiB & 6.4 MiB & 6.3 MiB & 6.3 MiB &   -9 KiB & Flat (Leak-Free) \\
\texttt{telegraf}                 &  0.3\% &   1.5\% & 45.0 MiB & 46.0 MiB & 44.2 MiB & 45.5 MiB &  +1.3 MiB & Flat \\
\texttt{mosquitto}                &  0.5\% &   3.5\% & 6.5 MiB & 11.9 MiB & 6.4 MiB & 6.4 MiB &  +13 KiB & Flat \\
\texttt{sim-control-panel}        &  0.1\% &   2.8\% & 26.4 MiB & 26.9 MiB & 26.6 MiB & 26.2 MiB & -421 KiB & Flat \\
\bottomrule
\end{tabular}
}
\end{table}

\paragraph{Subsystem Memory Analysis and Asymptotic Plateau.}
The extended multi-day observations provide conclusive empirical evidence of resource stabilization:
\begin{itemize}
    \item \textbf{Spark Structured Streaming JVM:} The resident heap occupancy of \texttt{spark-fleet-features} displayed a net downward trend over 38.9~hours (-7\% / -160~MiB), settling at 2.1~GiB. Observable major garbage collection cycles (such as at hour 30, where resident memory dropped to 1.81~GiB before returning to baseline) confirm that the stateless transformations accumulate no long-lived references.
    \item \textbf{MinIO S3 Storage:} MinIO resident memory entered a strictly bounded asymptotic plateau at $\sim0.72$~GiB (741~MiB), oscillating between 0.66~GiB and 0.75~GiB for more than 25 consecutive hours. Low-level cgroup telemetry (\texttt{memory.stat}) revealed that MinIO's application heap accounted for only 394.7~MiB; the remaining 355.2~MiB consisted of the Linux kernel's temporary directory entry cache (\texttt{slab\_reclaimable}), which the operating system evicts dynamically under memory pressure.
    \item \textbf{QuestDB Historian:} QuestDB maintained an essentially horizontal memory profile at $\sim1.55$~GiB (+6\% / +88~MiB drift over 39~hours). This minor increment corresponds to the operating system mapping virtual page headers (\texttt{mmap}) into memory as the database transitions across midnight to initialize new daily partition files (\texttt{PARTITION BY DAY}). Crucially, this memory footprint is strictly self-bounding and ceases expanding after the seventh operational day: with the 7-day retention policy active (\texttt{TTL 7 DAYS}), partitions older than one week are purged automatically, while the Linux page cache evicts dormant historical pages, locking QuestDB into a permanent 1.5--1.8~GiB steady-state operating band.
    \item \textbf{Redpanda Streaming Broker:} Memory reached its dynamic shard arena ceiling at $\sim3.0\text{--}3.1$~GiB (peaking at 3.06~GiB before settling at 3.03~GiB), after which further allocations from the host kernel ceased and internal buffer recycling stabilized the footprint permanently.
    \item \textbf{Grafana Observability Optimization:} During initial monitoring, raw timeseries queries lacking temporal aggregation transferred over 4.5~GB of unaggregated points across frequent polling cycles, causing browser UI sluggishness. Introducing server-side downsampling (\texttt{SAMPLE BY(\$\_\_interval)}) and adjusting the auto-refresh interval to 10~seconds eliminated this bottleneck, reducing network transfer payloads by 99.5\%, accelerating query execution to 2--50~ms, and keeping Grafana resident memory completely flat at 170~MiB.
    \item \textbf{Native Rust Simulator:} Maintained a completely flat, invariant resident footprint of 6.3~MiB with a negligible drift of only $-8.8$~KiB across the entire 39-hour run, demonstrating that strict compile-time ownership prevents memory leaks entirely.
\end{itemize}

Figure~\ref{fig:pipeline_memory_drift_phase2} charts the comparative memory trajectory over the 38.9-hour Phase~2 soak test.

\begin{figure}[htbp]
\centering
\begin{tikzpicture}
\begin{axis}[
    title={\textbf{Phase 2: Extended Multi-Day Memory Stabilization Trajectory (38.9 Hours)}},
    xlabel={Execution Time (Hours)},
    ylabel={RAM Consumption (GiB)},
    xmin=0, xmax=39,
    ymin=0, ymax=4.2,
    width=0.92\textwidth,
    height=7.2cm,
    grid=both,
    grid style={line width=.1pt, draw=gray!20},
    legend style={at={(0.02,0.96)}, anchor=north west, font=\scriptsize},
    thick
]
    % Redpanda
    \addplot[color=accentgreen, mark=diamond*] coordinates {
        (0, 1.68) (2, 1.76) (4, 1.85) (6, 1.93) (8, 2.01) (10, 2.10) (12, 2.18) (14, 2.26) 
        (16, 2.31) (18, 2.37) (20, 2.42) (22, 2.49) (24, 2.55) (26, 2.63) (28, 2.71) 
        (30, 2.80) (32, 2.86) (34, 2.93) (36, 3.00) (38, 3.06) (38.9, 3.03)
    };
    \addlegendentry{Redpanda Kafka (Seastar Arena Ceiling)}

    % Spark Fleet Features
    \addplot[color=primary, mark=square*] coordinates {
        (0, 2.35) (2, 2.36) (4, 2.37) (6, 2.38) (8, 2.38) (10, 2.39) (12, 2.40) (14, 2.41) 
        (16, 2.35) (18, 2.33) (20, 2.29) (22, 2.03) (24, 2.27) (26, 2.27) (28, 2.28) 
        (30, 1.81) (32, 2.28) (34, 2.16) (36, 2.27) (38, 2.27) (38.9, 2.10)
    };
    \addlegendentry{Spark Fleet Features (Stateless JVM GC)}

    % QuestDB
    \addplot[color=orange, mark=circle*] coordinates {
        (0, 1.39) (2, 1.47) (4, 1.47) (6, 1.48) (8, 1.49) (10, 1.50) (12, 1.51) (14, 1.48) 
        (16, 1.61) (18, 1.61) (20, 1.59) (22, 1.59) (24, 1.58) (26, 1.55) (28, 1.56) 
        (30, 1.57) (32, 1.58) (34, 1.58) (36, 1.58) (38, 1.55) (38.9, 1.55)
    };
    \addlegendentry{QuestDB (Columnar mmap Cache)}

    % MinIO (Hardened)
    \addplot[color=secondary, mark=triangle*] coordinates {
        (0, 0.51) (2, 0.53) (4, 0.56) (6, 0.59) (8, 0.63) (10, 0.65) (12, 0.70) (14, 0.72) 
        (16, 0.67) (18, 0.67) (20, 0.65) (22, 0.70) (24, 0.66) (26, 0.70) (28, 0.72) 
        (30, 0.75) (32, 0.74) (34, 0.75) (36, 0.73) (38, 0.75) (38.9, 0.72)
    };
    \addlegendentry{MinIO S3 (Bounded Plateau)}
\end{axis}
\end{tikzpicture}
\caption{Evolution of memory footprint over the extended 38.9-hour multi-day soak test. MinIO, Spark, and QuestDB exhibited horizontal stabilization, while Redpanda reached its pre-allocated Seastar arena ceiling, proving zero unbounded resource growth.}
\label{fig:pipeline_memory_drift_phase2}
\end{figure}

%======================================================================
\subsection{Consolidated Benchmark Summary and Conclusions}
\label{subsec:eval_summary}

Table~\ref{tab:pipeline_consolidated_summary} consolidates the core operational metrics across the baseline run and the hardened multi-day validation run, providing an executive comparison of the platform before and after architectural hardening.

\begin{table}[htbp]
\centering
\small
\caption{Consolidated Platform Metrics: Baseline Run vs. Hardened Multi-Day Soak Campaign}
\label{tab:pipeline_consolidated_summary}
\vspace{0.2cm}
\begin{tabular}{l r r l}
\toprule
\textbf{Evaluation Metric} & \textbf{Phase 1: Baseline (Day 1)} & \textbf{Phase 2: Hardened (Days 2--3)} & \textbf{Operational Impact} \\
\midrule
Metric Audit Duration      & 14.30 hours (1,641 snaps)  & \textbf{38.86 hours (4,450 snaps)}  & $+172\%$ sustained soak duration \\
Continuous Cluster Uptime  & 14.30 hours                & \textbf{62.35 hours (2.6 days)}    & 100\% availability (0 restarts) \\
Ingested Events (QuestDB)  & 233,412 rows               & \textbf{1,026,060 rows}            & $> 1$ Million rows persisted \\
Streaming Consumer Lag     & 0.0 (zero backlog)         & \textbf{0.0 (zero backlog)}        & Hard real-time stream parity \\
Total Lakehouse S3 Storage & 128.0 GiB (in 14.3h)       & \textbf{70.0 MiB (in 38.9h)}       & \textbf{-99.95\% storage reduction} \\
Retained Iceberg Files     & 61,757 files               & \textbf{1,073 files}               & Bound \& compacted footprint \\
Spark Batch Commit Time    & 2.5s -- 5.0s               & \textbf{0.18s}                     & \textbf{15$\times$ faster commits} \\
MinIO Resident Memory      & 2.07 GiB (+406\% growth)   & \textbf{0.72 GiB (stable plateau)} & Memory leak eliminated \\
Automated Lakehouse Cycles & 0 (inoperable manual call) & \textbf{87 successful cycles}      & 26.2s avg compaction cycle \\
\bottomrule
\end{tabular}
\end{table}

\paragraph{Concluding Assessment.}
As demonstrated across the full three-day continuous testing campaign, the telemetry data platform is robust, leak-free, and suitable for sustained, mission-critical industrial deployments. Automated metadata compaction and intelligent maintenance permanently resolved the $O(N^2)$ storage bloat discovered during the initial baseline run, while maintaining sub-second processing latencies across all components. By successfully ingesting, transforming, and persisting over one million high-frequency telemetry messages under strict real-time guarantees, the platform definitively proves its architectural readiness for predictive maintenance and anomaly detection of industrial mobile robots.


